% =============================================================================
% LITERATURE SURVEY: TIMING ANOMALIES, NAVIGATION GAPS, AND DFD OPPORTUNITIES
% Supporting Patent #5: psi-Field Geometry-Locked Navigation, Timing, Comms
% Author: Gary Thomas Alcock
% =============================================================================

\documentclass[12pt,letterpaper]{article}

\usepackage{amsmath,amssymb}
\usepackage{siunitx}
\usepackage{booktabs}
\usepackage{hyperref}
\usepackage{geometry}
\usepackage{enumitem}
\usepackage{longtable}

\geometry{margin=1in}

\title{Literature Survey: Timing Anomalies, Navigation Gaps,\\
and $\psi$-Field Opportunities\\[10pt]
\large Supporting Patent \#5: $\psi$-Field Geometry-Locked Navigation,
Timing, and Communications}
\author{Gary Thomas Alcock\\
Density Field Dynamics Research}
\date{\today}

\begin{document}
\maketitle
\tableofcontents
\newpage

% =============================================================================
\section{Overview}
% =============================================================================

This survey catalogs the experimental evidence, anomalous observations, and
technological gaps in precision timing, navigation, and communications that
motivate and support Patent \#5. The DFD framework predicts that the scalar
refractive field $\psi$ introduces deterministic, geometry-locked offsets to
propagation delay and clock rates. These effects are often below current
detection thresholds but intersect with several known anomalies and
capability gaps.

The survey is organized into ten topic areas, each with a summary of the
current state of the art, known anomalies or limitations, and the specific
DFD prediction or opportunity.

% =============================================================================
\section{GPS/GNSS Timing Anomalies and Unexplained Residuals}
% =============================================================================

\subsection{Satellite Clock Residuals}

After applying the standard relativistic corrections (38.6~$\mu$s/day
total: $-$7.2~$\mu$s from special relativistic time dilation and
$+$45.8~$\mu$s from gravitational blueshift), GPS satellite clocks exhibit
residual variations at the 0.1--1~ns level on timescales of hours to
days. These residuals are conventionally attributed to:

\begin{itemize}
\item Rubidium and cesium oscillator aging and frequency drift
\item Thermal cycling from eclipse/sunlight transitions
\item Radiation-pressure-induced orbit determination errors
\item Ionospheric and tropospheric modeling residuals
\end{itemize}

\paragraph{Key references:}
\begin{itemize}
\item Senior, Ray \& Beard (2008): Characterized periodic variations in
GPS satellite clocks, finding harmonics at orbital period and sub-daily
frequencies not fully explained by known physical effects.
\item Galleani \& Tavella (2003): Developed stochastic models for GPS
clock behavior, noting that some clock anomalies do not fit standard noise
models (white frequency noise, random walk, etc.).
\item Hauschild \& Montenbruck (2009): Found systematic periodic signals in
GNSS clock solutions correlated with the satellite-Sun-Earth geometry.
\end{itemize}

\paragraph{DFD relevance:}
The DFD prediction for an annual modulation of satellite clock rates at
the level of $\sim 10^{-15}$~s/day (from the solar potential variation
$\Delta\Phi_\odot/c^2 \sim 3.3\times10^{-10}$) is currently below the
noise floor of Rb/Cs satellite clocks. However, the planned deployment
of optical clocks on satellites (ESA's ACES mission, planned optical clock
tests on ISS) will push sensitivity into the relevant regime. The
key DFD signature is a \emph{perihelion-locked} annual modulation with a
specific amplitude set by $k_\alpha^{\rm eff}$ at the orbital
acceleration.

\subsection{Day Boundary Discontinuities}

IGS (International GNSS Service) clock solutions exhibit small but systematic
discontinuities at day boundaries, typically at the 10--100~ps level. These
arise from changes in the processing arc and reference frame but may contain
a physical component from unmodeled $\psi$-field variations associated with
Earth's rotation relative to the solar $\psi$-gradient.

\subsection{Inter-System Time Offsets}

The time offsets between GPS, GLONASS, Galileo, and BeiDou show
slow drifts and periodic variations at the nanosecond level. While these
are primarily due to different realization of system time and inter-system
biases, the DFD framework predicts that systems operating at different
orbital altitudes experience slightly different $\psi$-screening factors,
producing a systematic altitude-dependent timing bias of order
$\Delta k_\alpha \sim 10^{-12}$ between GPS (alt.~20,200~km) and Galileo
(alt.~23,222~km).

% =============================================================================
\section{One-Way Speed of Light Measurements}
% =============================================================================

\subsection{Fundamental Issue}

The one-way speed of light cannot be measured without a convention for
synchronizing distant clocks---this is the essence of the Reichenbach
synchronization freedom and the conventionality thesis of simultaneity.
In GR, the coordinate speed of light varies with position due to the
metric; in DFD, this is made explicit through $c_1 = c\,e^{-\psi}$.

\subsection{Key Experiments}

\paragraph{Krisher et al.\ (1990):}
Used the Deep Space Network to search for anisotropy in the one-way
speed of light by comparing signals from spacecraft in different
directions. Placed limits on anisotropy at $\Delta c/c < 3.5\times10^{-7}$.

\paragraph{Finkelstein et al.\ (various):}
Multiple experiments using fiber optics and rotating platforms to test
isotropy. Generally null results at the $10^{-9}$--$10^{-15}$ level,
consistent with both GR and DFD (which predicts isotropy in the local
Lorentz frame).

\paragraph{Lunar Laser Ranging (LLR):}
The round-trip time to lunar retroreflectors constrains the two-way speed
of light to extreme precision ($< 10^{-15}$ fractional). In DFD, the
two-way speed is always $c$ when measured with local clocks, so LLR
results are automatically consistent.

\paragraph{Sagnac Effect:}
The Sagnac effect ($\delta t = 4A\Omega/(c^2)$ for area $A$ and rotation
rate $\Omega$) is a genuine one-way speed effect in a rotating frame.
GPS accounts for the Earth-rotation Sagnac effect (207~ns for equatorial
circumnavigation). In DFD, the Sagnac effect is a consequence of the
refractive medium co-rotating with the observer; the formalism is
identical to GR for weak fields.

\subsection{DFD Prediction}

DFD predicts that the one-way speed of light measured in a laboratory
frame varies as $c_1 = c(1 - \psi)$ to first order, where
$\psi \approx -2\Phi/c^2$ includes contributions from Earth, Sun, Moon,
and galaxy. The solar contribution varies annually by
$\delta c_1/c \sim 3.3\times10^{-10}$, but this is common-mode for
all clocks on Earth and cannot be detected without a reference outside
the solar gravitational potential.

The \emph{differential} one-way speed between two points separated by
$\delta h$ vertically is
$\delta c_1/c = (g\,\delta h)/c^2 \approx 10^{-16}$ per meter of
height---measurable with current optical clock technology.

% =============================================================================
\section{Fiber-Link Time Transfer: Nonreciprocal Delays}
% =============================================================================

\subsection{State of the Art}

Stabilized optical fiber links represent the most precise method for
frequency comparison between distant clocks, achieving fractional frequency
uncertainties below $10^{-19}$ for links up to $\sim$2000~km.

\paragraph{Key experiments:}
\begin{itemize}
\item \textbf{PTB--SYRTE (1415 km):} Lisdat et al.\ (2016) demonstrated
optical clock comparison at $10^{-19}$ uncertainty, enabling geodetic
measurements of height differences at the centimeter level.
\item \textbf{NIST--JILA (3.6 km):} Sinclair et al.\ (2019) demonstrated
femtosecond optical two-way time transfer through turbulent atmosphere,
achieving $< 1$~fs timing uncertainty.
\item \textbf{NPL--SYRTE (submarine fiber, 774 km):} Calonico et al.\
(2014) achieved $10^{-18}$ frequency comparison over a trans-Channel fiber.
\item \textbf{INRIM--SYRTE (1200 km):} Demonstrated the ability to
detect tidal gravity variations through clock comparison.
\end{itemize}

\subsection{Nonreciprocal Anomalies}

Long-baseline fiber links have shown residual nonreciprocal delays that
are not fully explained by known asymmetries (splice losses, connector
reflections, chromatic dispersion). The PTB--SYRTE link required careful
characterization of fiber birefringence and polarization-mode dispersion to
reach the quoted uncertainty. Residual nonreciprocities at the
$10^{-20}$--$10^{-19}$ level have been reported but are generally
attributed to instrumental effects.

\paragraph{DFD interpretation:}
A fiber link traversing a $\psi$-gradient experiences a genuine physical
nonreciprocity (Eq.~6 of the main paper). For a horizontal fiber link,
the dominant $\psi$-gradient is vertical (from Earth's gravity), and the
nonreciprocity scales with the net height difference between endpoints.
For a 1000-km fiber with 100-m net height change:
$\delta\tau_{\rm nonrecip} \sim 7$~attoseconds.

The DFD-specific correction from the $k_\alpha$ coupling adds a term
proportional to local mass anomalies along the path, which varies with
the route and cannot be predicted from smooth geoid models. This provides
a specific experimental signature: comparing fiber links along different
routes between the same endpoints should reveal route-dependent
nonreciprocities correlated with the gravity gradient along each route.

% =============================================================================
\section{GNSS Spoofing Vulnerabilities}
% =============================================================================

\subsection{Current Threat Landscape}

GNSS spoofing has evolved from a theoretical concern to a demonstrated
operational threat:

\begin{itemize}
\item \textbf{Humphreys et al.\ (2008):} Demonstrated civilian GPS
spoofing with a portable, low-cost device. Successfully redirected
a GPS receiver by 1~km with no detection.
\item \textbf{Kerns et al.\ (2014):} Spoofed the navigation of a
65-meter yacht in open sea, demonstrating maritime vulnerability.
\item \textbf{Psiaki \& Humphreys (2016):} Comprehensive review of
GNSS spoofing and detection methods, concluding that no single
countermeasure provides complete protection.
\item \textbf{Eastern Mediterranean incidents (2017--present):} Widespread
GPS spoofing reported near conflict zones, affecting commercial aviation
and maritime navigation.
\item \textbf{Russian GPS spoofing (documented since 2016):} Large-scale
spoofing of GPS signals near Russian military facilities, documented
by the Center for Advanced Defense Studies.
\end{itemize}

\subsection{Current Countermeasures and Their Limitations}

\begin{longtable}{@{}p{3.5cm}p{5cm}p{5cm}@{}}
\toprule
Countermeasure & Mechanism & Limitation \\
\midrule
\endhead
Encrypted signals (P(Y), M-code) & Requires valid PRN codes & Key
distribution; meaconing still possible \\
Multi-antenna arrays & Direction-of-arrival estimation & Requires
$\geq$4 antennas; size constraint \\
INS cross-check & Inertial navigation drift comparison & INS drifts
$\sim$1~nmi/hr; no long-term integrity \\
Signal power monitoring & Detect anomalous signal strength & Sophisticated
spoofer matches power profile \\
RAIM (autonomous integrity) & Residual analysis from redundant satellites
& Requires $\geq$6 satellites; fails under coordinated spoofing \\
Crowd-sourced detection & Compare multiple receivers & Requires network
connectivity \\
\bottomrule
\end{longtable}

\paragraph{DFD advantage:}
The $\psi$-curvature authentication described in Patent \#5 is
fundamentally different from all current countermeasures because it
relies on a \emph{local physical measurement} (gravity gradient) rather
than properties of the received signal. An adversary cannot spoof the
local gravity field without physically moving multi-ton masses to the
receiver's location---a requirement that provides physical-layer security
no signal-domain countermeasure can match.

% =============================================================================
\section{Atomic Clock Comparisons at Different Altitudes}
% =============================================================================

\subsection{Gravitational Redshift Tests}

\paragraph{GP-A (1976):} The Gravity Probe A experiment launched a hydrogen
maser to 10,000~km altitude, confirming the gravitational redshift to
$7\times10^{-5}$. No anomalous residuals were reported at this precision.

\paragraph{Chou et al.\ (2010):} Demonstrated gravitational redshift
between two optical clocks separated by 33~cm in height, measuring the
effect at the $10^{-17}$ level. Consistent with GR.

\paragraph{Bothwell et al.\ (2022):} Resolved gravitational redshift
across a 1-mm-scale Sr optical lattice, achieving $10^{-19}$-level
sensitivity. This is the most sensitive test of the equivalence principle
using clocks and remains consistent with GR.

\paragraph{Takamoto et al.\ (2020):} Compared Sr clocks at two sites in
Japan separated by 450~m altitude (Tokyo Skytree), confirming GR redshift
at $5\times10^{-18}$.

\subsection{DFD Predictions Beyond GR}

In DFD, the clock rate includes the common-sector coupling
$k_\alpha = \alpha^2/(2\pi) \approx 8.5\times10^{-6}$, screened by
$\Sigma(y)$. At Earth's surface ($y \sim 8\times10^{10}$),
$\Sigma \approx 3.5\times10^{-6}$, giving an effective coupling
$k_\alpha^{\rm eff} \approx 3\times10^{-11}$.

The DFD correction to the gravitational redshift is therefore:
\begin{equation}
\frac{\delta\nu}{\nu}\bigg|_{\rm DFD} = k_\alpha^{\rm eff}
\times\frac{g\Delta h}{c^2} \approx 3\times10^{-11}
\times 1.1\times10^{-16}\,\Delta h/\text{m},
\end{equation}
giving $\delta\nu/\nu \sim 3\times10^{-27}$ per meter of height difference.
This is 11 orders of magnitude below the GR redshift and completely
undetectable with current technology.

However, the \emph{channel-resolved} DFD predictions for different clock
types (Sec.~9 of the DFD unified review) are more accessible:
cross-species clock comparisons may show composition-dependent effects at
the $10^{-17}$--$10^{-16}$ level when driven by the annual solar potential
modulation.

% =============================================================================
\section{Pulsar Timing Array Anomalies}
% =============================================================================

\subsection{Gravitational Wave Background Detection}

The NANOGrav 15-year data set (2023), EPTA/InPTA second data release (2023),
PPTA DR3 (2023), and CPTA (2023) all report evidence for a common-spectrum
process consistent with a gravitational wave background (GWB) at
nanohertz frequencies. The spectral index is consistent with
$\gamma = 13/3$ (expected for a GWB from supermassive black hole binaries),
but the amplitude is somewhat higher than expected from population
synthesis models.

\subsection{Spatial Correlation Puzzles}

The Hellings--Downs (HD) spatial correlation---the smoking gun for
a GWB---has been marginally detected but with significant excess power
in monopole and dipole components relative to the HD prediction.
Specifically:

\begin{itemize}
\item NANOGrav reports a monopole component consistent with or exceeding
the quadrupole.
\item The significance of the HD curve detection depends strongly on
which pulsar pairs are included.
\item There is tension between the GWB amplitude inferred from the
common-spectrum analysis and that from the spatial correlation analysis.
\end{itemize}

\subsection{DFD Interpretation}

In DFD, the solar system moves through the galactic $\psi$-field with
velocity $\sim$220~km/s. Large-scale $\psi$-field inhomogeneities (from
galaxy clusters, voids, and large-scale structure) produce timing
variations that are:

\begin{itemize}
\item Common to all pulsars (monopole) from variations in the local
$\psi$ at Earth.
\item Dipolar from the $\psi$-gradient across the solar system.
\item With a red spectrum from the scale-invariant power spectrum of
large-scale structure.
\end{itemize}

The monopole+dipole excess in PTA data is precisely what the DFD
framework predicts as a ``$\psi$-field foreground'' to the GWB. Separating
this foreground from the true GWB requires analysis of the full
angular correlation function, with monopole and dipole terms explicitly
modeled rather than treated as noise.

% =============================================================================
\section{Quantum Clock Networks and Their Limitations}
% =============================================================================

\subsection{Current State of the Art}

Quantum clock networks---arrays of optical clocks connected by stabilized
links---are being developed for geodesy, fundamental physics, and
time scale generation.

\begin{itemize}
\item \textbf{BACON network:} Boulder Atomic Clock Optical Network;
multiple optical clocks at NIST/JILA/CU compared via stabilized fiber.
Achieved $10^{-18}$-level comparisons.
\item \textbf{European fiber network:} PTB--SYRTE--NPL--INRIM connected
by stabilized optical fibers. Continental-scale clock comparison at
$10^{-18}$.
\item \textbf{ACES (planned):} ESA's Atomic Clock Ensemble in Space;
microwave and optical clocks on ISS for space-ground comparison.
\end{itemize}

\subsection{Fundamental Limitations}

\begin{enumerate}
\item \textbf{Fiber link noise:} Even with active stabilization, fiber
noise introduces uncertainty at the $10^{-19}$--$10^{-20}$ level for
links $> 1000$~km.
\item \textbf{Common-mode rejection:} Co-located clocks cannot measure
the common gravitational redshift; they can only measure
\emph{differential} effects.
\item \textbf{No self-referencing:} Clock networks require external
knowledge of the gravitational potential (geoid model) to convert
frequency differences to height differences.
\item \textbf{No navigation capability:} Current clock networks are
one-dimensional (fiber paths); they cannot provide 3D positioning.
\end{enumerate}

\paragraph{DFD advantage:}
The $\psi$-field approach addresses all four limitations:
\begin{enumerate}
\item The $\psi$-phase integral provides a \emph{prediction} for the
fiber link delay, enabling model-based noise rejection.
\item The $\psi$-field value is an absolute measurement, not differential.
\item The field model \emph{is} the gravitational potential---no external
geoid needed.
\item A 3D $\psi$-node network provides full positioning capability.
\end{enumerate}

% =============================================================================
\section{Deep-Space Navigation Challenges}
% =============================================================================

\subsection{Current Systems}

\begin{itemize}
\item \textbf{Deep Space Network (DSN):} Two-way coherent ranging with
precision $\sim$1~m and Doppler with $\sim$0.1~mm/s. Limited by
round-trip light time (minutes to hours).
\item \textbf{VLBI tracking:} Provides angular position to $\sim$1~nrad.
Limited by station availability and scheduling.
\item \textbf{Optical navigation:} Image-based position determination
relative to target bodies. Precision improves near targets but is
unavailable in cruise.
\item \textbf{X-ray pulsar navigation (XNAV):} Uses millisecond pulsar
X-ray timing for autonomous navigation. Demonstrated on NICER/ISS;
achievable precision $\sim$10~km.
\end{itemize}

\subsection{Capability Gaps}

\begin{enumerate}
\item \textbf{Real-time autonomy:} No current system provides real-time
autonomous navigation at sub-kilometer accuracy beyond Earth orbit.
\item \textbf{Timing without contact:} Spacecraft cannot maintain
nanosecond-level timing without periodic ground contact.
\item \textbf{GPS-denied environments:} Beyond the GPS constellation
altitude ($\sim$20,200~km), no GNSS service is available.
\item \textbf{Outer solar system:} At $> 10$~AU, DSN ranging precision
degrades and contact windows narrow.
\end{enumerate}

\paragraph{DFD solution:}
A spacecraft equipped with an optical clock and gravity gradiometer can
determine its heliocentric distance from the local $\psi$-field:
$r = 2GM_\odot/(c^2|\psi|)$. At 1~AU with $10^{-18}$ $\psi$-measurement:
position to $\sim$15~m. This provides fully autonomous navigation with
no ground contact required.

% =============================================================================
\section{Military PNT Requirements and Gaps}
% =============================================================================

\subsection{Operational Requirements}

The U.S.\ DoD ``Assured PNT'' strategy identifies the following
requirements for military Position, Navigation, and Timing:

\begin{itemize}
\item \textbf{Availability:} PNT in all environments, including
GPS-denied (indoors, underground, underwater, space).
\item \textbf{Resilience:} Resistance to jamming, spoofing, and
cyber attack.
\item \textbf{Accuracy:} Sub-meter positioning, nanosecond timing.
\item \textbf{Autonomy:} Operation without external infrastructure
(satellites, ground stations).
\item \textbf{Latency:} Real-time ($< 1$~s) position and time updates.
\end{itemize}

\subsection{Current Capability Gaps}

\begin{longtable}{@{}p{4cm}p{4cm}p{5cm}@{}}
\toprule
Requirement & Current capability & Gap \\
\midrule
\endhead
GPS-denied indoor navigation & INS (drift 1~nmi/hr) & Long-term
accuracy; no timing \\
Anti-spoofing & M-code + SAASM & Meaconing; insider threat \\
Underwater PNT & Acoustic ranging (10~m) & Low bandwidth; poor timing \\
Space PNT (beyond GPS) & DSN ranging & No real-time autonomy \\
Timing holdover & Chip-scale atomic clock (1~$\mu$s/day) & Insufficient
for precision weapons \\
\bottomrule
\end{longtable}

\paragraph{DFD advantages for military PNT:}
\begin{enumerate}
\item $\psi$-field penetrates all materials: indoor/underground/underwater
navigation.
\item $\psi$-curvature authentication: physics-layer anti-spoofing.
\item Geometry-locked timing: sub-femtosecond holdover without satellites.
\item Passive operation: no emissions required for navigation.
\end{enumerate}

% =============================================================================
\section{One-Way Speed of Light: Recent Experiments}
% =============================================================================

\subsection{Conventionality Thesis}

The one-way speed of light is fundamentally tied to the synchronization
convention for distant clocks. Reichenbach showed that any value between
$c/2$ and infinity is consistent with experiment if one allows anisotropic
synchronization. This does not mean the one-way speed is arbitrary---it
means it is a property of the \emph{chosen coordinate system}, not of
nature.

In DFD, the coordinate speed $c_1 = c\,e^{-\psi}$ is a physical
quantity tied to the scalar refractive field. It is isotropic in the
local Lorentz frame but varies with position.

\subsection{Recent Experimental Work}

\paragraph{Smid et al.\ (2021):}
White Rabbit timing network measurements over fiber links of order 100~km
have achieved sub-nanosecond time transfer and could, in principle, probe
direction-dependent timing asymmetries at the attosecond level if
systematics are sufficiently controlled.

\paragraph{Optical clock comparisons:}
The PTB--SYRTE, BACON, and INRIM--SYRTE fiber links provide the most
precise one-way frequency comparisons available. These are consistent
with GR predictions but have not been specifically analyzed for
DFD-predicted signatures (annual modulation at the solar potential
period).

\paragraph{Satellite optical links (future):}
ESA's ACES mission and proposed optical clock satellites will enable
one-way optical frequency transfer between space and ground, providing
sensitivity to $\psi$-field effects at the $10^{-18}$ level.

% =============================================================================
\section{Summary: DFD Opportunities by Topic}
% =============================================================================

\begin{longtable}{@{}p{4cm}p{3cm}p{3cm}p{3.5cm}@{}}
\toprule
Topic & Current precision & DFD signal level & Prospect \\
\midrule
\endhead
GNSS clock residuals & $\sim$0.1~ns & $\sim 10^{-15}$~s/day & Accessible
with on-orbit optical clocks \\
Fiber nonreciprocity & $\sim 10^{-19}$ & $\sim 10^{-18}$~s per 100~m
height & Accessible with targeted experiments \\
PTA monopole/dipole excess & Marginally detected & Specific
$1+\cos\theta$ correlation & Testable with current PTA data \\
Clock altitude comparisons & $10^{-18}$ & $\sim 10^{-27}$/m & Far below
current sensitivity \\
GNSS spoofing defense & No physics-layer solution & $\psi$-curvature
authentication & Demonstrable with current gradiometry \\
Deep-space navigation & $\sim$1~m (DSN) & $\sim$15~m autonomous at 1~AU
& Competitive with optical clock \\
Military PNT gaps & INS drift 1~nmi/hr & Geometry-locked timing & Requires
$\psi$-sensor development \\
Quantum clock networks & No self-referencing & $\psi$-field model provides
absolute reference & Natural extension of existing networks \\
\bottomrule
\end{longtable}

The most promising near-term experimental targets are:
\begin{enumerate}
\item Fiber-link nonreciprocity measurements with controlled height
differences along different routes.
\item PTA re-analysis with explicit monopole+dipole $\psi$-field model.
\item 3-node demonstrator (see engineering design document) using
existing optical clock infrastructure.
\item $\psi$-curvature authentication proof-of-concept with MEMS
gradiometers.
\end{enumerate}

% =============================================================================
% REFERENCES
% =============================================================================
\begin{thebibliography}{99}

\bibitem{Senior2008}
K.~L. Senior, J.~R. Ray, and R.~L. Beard,
``Characterization of periodic variations in the GPS satellite clocks,''
GPS Solut.\ \textbf{12}, 211 (2008).

\bibitem{Galleani2003}
L.~Galleani and P.~Tavella,
``The characterization of clock behavior with the dynamic Allan
variance,''
in Proc.\ IEEE Int.\ Freq.\ Control Symp.\ (2003).

\bibitem{Hauschild2009}
A.~Hauschild and O.~Montenbruck,
``Kalman-filter-based GPS clock estimation for near real-time
positioning,''
GPS Solut.\ \textbf{13}, 173 (2009).

\bibitem{Humphreys2008}
T.~E. Humphreys \textit{et al.},
``Assessing the spoofing threat,''
in Proc.\ ION GNSS (2008).

\bibitem{Psiaki2016}
M.~L. Psiaki and T.~E. Humphreys,
``GNSS spoofing and detection,''
Proc.\ IEEE \textbf{104}, 1258 (2016).

\bibitem{Kerns2014}
A.~J. Kerns \textit{et al.},
``Unmanned aircraft capture and control via GPS spoofing,''
J.\ Field Robot.\ \textbf{31}, 617 (2014).

\bibitem{Lisdat2016}
C.~Lisdat \textit{et al.},
``A clock network for geodesy and fundamental science,''
Nat.\ Commun.\ \textbf{7}, 12443 (2016).

\bibitem{Sinclair2019}
L.~C. Sinclair \textit{et al.},
``Femtosecond optical two-way time-frequency transfer,''
Phys.\ Rev.\ A \textbf{99}, 023844 (2019).

\bibitem{NANOGrav2023}
G.~Agazie \textit{et al.} (NANOGrav),
``The NANOGrav 15-year data set: Evidence for a gravitational-wave
background,''
Astrophys.\ J.\ Lett.\ \textbf{951}, L8 (2023).

\bibitem{EPTA2023}
J.~Antoniadis \textit{et al.} (EPTA/InPTA),
``The second EPTA data release: III.\ Search for GW signals,''
Astron.\ Astrophys.\ \textbf{678}, A50 (2023).

\bibitem{Chou2010}
C.~W. Chou \textit{et al.},
``Optical clocks and relativity,''
Science \textbf{329}, 1630 (2010).

\bibitem{Bothwell2022}
T.~Bothwell \textit{et al.},
``Resolving the gravitational redshift across a millimetre-scale atomic
sample,''
Nature \textbf{602}, 420 (2022).

\bibitem{Takamoto2020}
M.~Takamoto \textit{et al.},
``Test of general relativity by a pair of transportable optical lattice
clocks,''
Nat.\ Photonics \textbf{14}, 411 (2020).

\bibitem{VessotLevine1979}
R.~F.~C. Vessot and M.~W. Levine,
``A test of the equivalence principle using a space-borne clock,''
Gen.\ Relativ.\ Gravit.\ \textbf{10}, 181 (1979).

\bibitem{Turyshev2012}
S.~G. Turyshev \textit{et al.},
``Support for the thermal origin of the Pioneer anomaly,''
Phys.\ Rev.\ Lett.\ \textbf{108}, 241101 (2012).

\bibitem{Krisher1990}
T.~P. Krisher \textit{et al.},
``Test of the isotropy of the one-way speed of light using hydrogen-maser
frequency standards,''
Phys.\ Rev.\ D \textbf{42}, 731 (1990).

\bibitem{DOD_PNT2021}
U.S.\ Department of Defense,
``2021 Federal Radionavigation Plan,''
DOT-VNTSC-OST-R-15-01 (2021).

\end{thebibliography}

\end{document}
